\documentclass[11pt]{article}
\usepackage[margin=25mm]{geometry}
\usepackage{fontspec}
\setmainfont{Latin Modern Roman}
\usepackage{amsmath,amssymb,mathtools}
\usepackage{unicode-math}
\setmathfont{Latin Modern Math}
\usepackage{xurl,hyperref}
\hypersetup{colorlinks=true,urlcolor=blue}
\setcounter{secnumdepth}{0}
\setlength{\parindent}{0pt}
\setlength{\parskip}{5pt}
\setlength{\emergencystretch}{3em}
\newcommand{\sref}[2]{\hyperlink{source:#1:#2}{[S#2]}}
\newcommand{\cataloguescope}[1]{\paragraph{Recorded target.}#1}
\begin{document}
\section{Problems Around Polynomials — Conjecture 4}\label{statement}

\subsection{Definitions and mathematical statement}

\paragraph{Shared notation.}$\mathbb N=\{1,2,\ldots\}$, and $\mathbb Z,\mathbb Q,\mathbb R,\mathbb C$ denote the integers, rationals, reals and complex numbers. Any other conventions are specified locally.


For positive integers $k,\ell$, let $\widetilde\sharp(2k,\ell)$ be the largest possible number of isolated real zeros of a polynomial $p\in\mathbb R[x_1,\ldots,x_\ell]$ of total degree exactly $2k$ admitting a representation
\[p=\sum_{j=1}^m q_j^2,\qquad m\geq1,\quad\deg q_j\leq k.\]
Such a polynomial is nonnegative on $\mathbb R^\ell$. A zero $x$ is isolated if some neighborhood of $x$ contains no other zero; distinct points are counted once.
\paragraph{Statement recorded in the supplied source.}
Is it true for every $k,\ell\geq1$ that
\[\widetilde\sharp(2k,\ell)=k^\ell?\] \sref{statement}{2}
\subsection{Short English statement}
Can a degree-2k sum-of-squares polynomial in any number of variables have more than k to that number of isolated real zeros?
\subsection{Sources}
\begin{description}
\item[\hypertarget{source:statement:1}{S1}] ULAM.AI and the UnsolvedMath Contributors, UnsolvedMath: A Curated Collection of Open Mathematics Problems, catalogue entry. CC BY 4.0. \url{https://huggingface.co/datasets/ulamai/UnsolvedMath}.
\item[\hypertarget{source:statement:2}{S2}] Boris Shapiro, Problems Around Polynomials -- The Good, The Bad and The Ugly, Arnold Mathematical Journal 1 (2015), pp.91--104, Conjecture 4 and its complete preceding definitions. Original question and full discussion. \url{https://arxiv.org/abs/1503.05295}.

\end{description}
\label{end:statement}
\end{document}
